\documentclass[11pt,twoside,openright]{book}
\usepackage[margin=1in]{geometry}
\usepackage{xcolor}
\usepackage{listings}
\usepackage{fancyhdr}
\usepackage{graphicx}
\usepackage{hyperref}
\usepackage{amsmath}
\usepackage{amssymb}
\usepackage{tikz}
\usepackage{array}
\usepackage{booktabs}
\usepackage{syntax}
\usepackage{setspace}

% Define colors for syntax highlighting
\definecolor{codeblue}{rgb}{0.0, 0.3, 0.6}
\definecolor{codegreen}{rgb}{0.0, 0.6, 0.0}
\definecolor{codegray}{rgb}{0.5, 0.5, 0.5}
\definecolor{codepurple}{rgb}{0.58, 0.0, 0.82}
\definecolor{codecrimson}{rgb}{0.86, 0.08, 0.24}
\definecolor{codelightgray}{rgb}{0.95, 0.95, 0.95}

% Configure listings
\lstdefinelanguage{ConstructLang}{
  keywords={let,fn,if,then,else,match,type,import,include,true,false},
  keywordstyle=\color{codeblue}\bfseries,
  ndkeywords={Int,Float,Bool,String,List},
  ndkeywordstyle=\color{codepurple},
  identifierstyle=\color{black},
  sensitive=true,
  comment=[l]{\#},
  commentstyle=\color{codegreen}\ttfamily,
  morecomment=[s]{/*}{*/},
  stringstyle=\color{codecrimson},
  breaklines=true,
  breakatwhitespace=false,
  tabsize=2,
  showstringspaces=false,
}

\lstset{
  language=ConstructLang,
  backgroundcolor=\color{codelightgray},
  basicstyle=\small\ttfamily,
  numbers=left,
  numberstyle=\tiny\color{codegray},
  numbersep=5pt,
  framexleftmargin=10pt,
  frame=tlrb,
  rulecolor=\color{codegray},
  captionpos=b,
  xleftmargin=15pt,
}

% Set up page style
\pagestyle{fancy}
\setlength{\headheight}{14pt}
\fancyhf{}
\fancyhead[LE,RO]{\thepage}
\fancyhead[RE]{\textit{\leftmark}}
\fancyhead[LO]{\textit{\rightmark}}
\renewcommand{\headrulewidth}{0.5pt}

% Title and author
\title{\Huge \textbf{The Construct Programming Language}\\
        \Large A Comprehensive Guide to Writing Functional Programs}
\author{\Large Charlie Santana}
\date{\Large October 28, 2025}

\hypersetup{
  colorlinks=true,
  linkcolor=codeblue,
  filecolor=magenta,
  urlcolor=codeblue,
}

\begin{document}

\frontmatter
\maketitle

\chapter*{Preface}

Welcome to \textbf{Construct}, a functional programming language that brings together the elegance of functional paradigms with practical usability. This book serves as a comprehensive guide to writing, understanding, and mastering the Construct language.

Construct is a statically-typed, functional language that compiles to LLVM IR and produces native binaries. It combines the power of languages like Haskell and ML with syntax that is approachable and expressive.

Whether you're new to functional programming or an experienced developer exploring new paradigms, this guide will take you through:

\begin{itemize}
    \item The fundamentals of Construct syntax
    \item Control flow and expressions
    \item Functions and lambdas
    \item Type system and type annotations
    \item Standard library and built-in functions
    \item Practical examples and patterns
    \item Advanced features and techniques
\end{itemize}

\section*{How to Use This Book}

This book is organized to support multiple reading approaches:

\begin{itemize}
    \item \textbf{Sequential}: Read from beginning to end for a complete introduction
    \item \textbf{By Topic}: Jump to chapters covering specific features
    \item \textbf{Reference}: Use the index to quickly find syntax and functions
    \item \textbf{Examples}: Study the code examples to learn through practice
\end{itemize}

\section*{Prerequisites}

This book assumes:

\begin{itemize}
    \item Basic familiarity with programming concepts (variables, functions, control flow)
    \item Comfort with reading code examples
    \item Access to a Construct compiler (installation instructions in Appendix A)
\end{itemize}

Experience with functional programming is helpful but not required.

\newpage
\tableofcontents

\mainmatter

\chapter{Introduction to Construct}

\section{What is Construct?}

Construct is a \textbf{statically-typed functional programming language} designed to be:

\begin{description}
    \item[Functional] Everything is an expression; immutability by default
    \item[Practical] Accessible syntax without sacrificing power
    \item[Safe] Static typing prevents entire classes of errors
    \item[Fast] Compiles to native code via LLVM
    \item[Expressive] Pattern matching, type inference, and first-class functions
\end{description}

\section{Design Philosophy}

Construct is built on these core principles:

\begin{enumerate}
    \item \textbf{Expressions over Statements}: Programs are composed of expressions that return values
    \item \textbf{Immutability by Default}: Values don't change; new values are created instead
    \item \textbf{Type Safety}: Static types catch errors at compile-time
    \item \textbf{Composition}: Functions are first-class and can be combined easily
    \item \textbf{Clarity}: Syntax should be readable and intention-revealing
\end{enumerate}

\section{A First Program}

Let's look at the simplest program in Construct:

\begin{lstlisting}[language=ConstructLang]
5 + 3
\end{lstlisting}

This program outputs the result: \texttt{8}

Here's something more interactive:

\begin{lstlisting}[language=ConstructLang]
let x: Int = 42
showln("The answer is ", x)
\end{lstlisting}

This program:
\begin{enumerate}
    \item Binds the value \texttt{42} to the name \texttt{x}
    \item Prints ``The answer is 42'' followed by a newline
\end{enumerate}

\section{Running Construct Programs}

To compile and run a Construct program:

\begin{lstlisting}[language=bash]
# Compile to LLVM IR
construct myprogram.ct

# Compile to native binary
construct myprogram.ct -c
clang++ -o myprogram myprogram.o ./libconstruct_stdlib.so -lm

# Run the binary
./myprogram
\end{lstlisting}

\chapter{Syntax Fundamentals}

\section{Tokens and Literals}

\subsection{Integers}

Integer literals are whole numbers:

\begin{lstlisting}[language=ConstructLang]
42
0
-100
1000000
\end{lstlisting}

\subsection{Floating-Point Numbers}

Floating-point literals include a decimal point:

\begin{lstlisting}[language=ConstructLang]
3.14
-2.5
0.0
1e6
\end{lstlisting}

\subsection{Strings}

String literals are enclosed in double quotes:

\begin{lstlisting}[language=ConstructLang]
"Hello, Construct!"
"Line 1\nLine 2"
""
\end{lstlisting}

\subsection{Booleans}

Boolean values are \texttt{true} or \texttt{false}:

\begin{lstlisting}[language=ConstructLang]
true
false
\end{lstlisting}

\section{Comments}

Single-line comments start with \texttt{\#}:

\begin{lstlisting}[language=ConstructLang]
# This is a comment
let x = 5  # Inline comment
\end{lstlisting}

Multi-line comments are enclosed in \texttt{/* */}:

\begin{lstlisting}[language=ConstructLang]
/*
  This is a multi-line comment
  It can span many lines
*/
let y = 10
\end{lstlisting}

\section{Indentation}

Construct uses Python-style indentation to define code blocks. Indentation is \textbf{significant}:

\begin{lstlisting}[language=ConstructLang]
let x = 1
let y = 2
showln(x + y)
\end{lstlisting}

Indentation increases inside function bodies, match expressions, and other blocks. Improper indentation results in syntax errors.

\section{Math Functions}

\begin{center}
\begin{tabular}{|c|c|l|}
\hline
\textbf{Operator} & \textbf{Name} & \textbf{Example} \\
\hline
\texttt{+} & Addition & \texttt{5 + 3} \\
\texttt{-} & Subtraction & \texttt{10 - 4} \\
\texttt{*} & Multiplication & \texttt{6 * 7} \\
\texttt{/} & Division & \texttt{20 / 4} \\
\texttt{\%} & Modulo & \texttt{17 \% 5} \\
\texttt{**} & Power & \texttt{2 ** 8} \\
\hline
\end{tabular}
\end{center}

\subsection{Comparison Operators}

\begin{center}
\begin{tabular}{|c|c|l|}
\hline
\textbf{Operator} & \textbf{Meaning} & \textbf{Example} \\
\hline
\texttt{==} & Equal to & \texttt{5 == 5} \\
\texttt{!=} & Not equal to & \texttt{5 != 3} \\
\texttt{<} & Less than & \texttt{3 < 5} \\
\texttt{<=} & Less or equal & \texttt{5 <= 5} \\
\texttt{>} & Greater than & \texttt{5 > 3} \\
\texttt{>=} & Greater or equal & \texttt{5 >= 5} \\
\hline
\end{tabular}
\end{center}

\section{Logical Operators}

\begin{center}
\begin{tabular}{|c|c|l|}
\hline
\textbf{Operator} & \textbf{Meaning} & \textbf{Example} \\
\hline
\texttt{\&\&} & Logical AND & \texttt{true \&\& false} \\
\texttt{||} & Logical OR & \texttt{true || false} \\
\texttt{!} & Logical NOT & \texttt{!true} \\
\hline
\end{tabular}
\end{center}

\section{String Concatenation}

The \texttt{+} operator concatenates strings:

\begin{lstlisting}[language=ConstructLang]
"Hello" + " " + "World"
"Numbers: " + toStr(42)
\end{lstlisting}

\section{Operator Precedence}

Operators are evaluated in this order (highest to lowest precedence):

\begin{enumerate}
    \item Power (\texttt{**})
    \item Multiplication, Division, Modulo (\texttt{*}, \texttt{/}, \texttt{\%})
    \item Addition, Subtraction (\texttt{+}, \texttt{-})
    \item Comparison (\texttt{<}, \texttt{<=}, \texttt{>}, \texttt{>=})
    \item Equality (\texttt{==}, \texttt{!=})
    \item Logical AND (\texttt{\&\&})
    \item Logical OR (\texttt{||})
\end{enumerate}

Use parentheses to override precedence:

\begin{lstlisting}[language=ConstructLang]
2 + 3 * 4      # Evaluates to 14
(2 + 3) * 4    # Evaluates to 20
\end{lstlisting}

\chapter{Variables and Types}

\section{Variable Binding}

Variables are bound using \texttt{let}:

\begin{lstlisting}[language=ConstructLang]
let x = 5
let name = "Alice"
let pi = 3.14159
\end{lstlisting}

Variables in Construct are \textbf{immutable}. Once bound, their value cannot change:

\begin{lstlisting}[language=ConstructLang]
let x = 10
let x = 20  # Error: x is already defined
\end{lstlisting}

\section{Type Annotations}

You can explicitly annotate types. Types include:

\begin{itemize}
    \item \texttt{Int}: 32-bit integers
    \item \texttt{Float}: 64-bit floating-point numbers
    \item \texttt{Bool}: Boolean values
    \item \texttt{String}: Text strings
    \item Function types: \texttt{Int -> Int}
\end{itemize}

\subsection{Basic Type Annotations}

\begin{lstlisting}[language=ConstructLang]
let x: Int = 42
let y: Float = 3.14
let name: String = "Bob"
let flag: Bool = true
\end{lstlisting}

\subsection{Function Type Annotations}

Function types are written with arrows:

\begin{lstlisting}[language=ConstructLang]
# A function that takes an Int and returns an Int
double: Int -> Int = fn (x: Int) -> x * 2

# A function that takes two Ints and returns an Int
add: Int -> Int -> Int = fn (a: Int) -> fn (b: Int) -> a + b
\end{lstlisting}

\section{Type Inference}

Construct can infer types from context. These are equivalent:

\begin{lstlisting}[language=ConstructLang]
let x = 5           # Type inferred as Int
let x: Int = 5      # Type explicitly annotated
\end{lstlisting}

Type inference works through:

\begin{itemize}
    \item Literal types (\texttt{5} is Int, \texttt{3.14} is Float)
    \item Function return types
    \item Context from operations
\end{itemize}

However, explicit type annotations are recommended for clarity and error detection.

\chapter{Control Flow}

\section{If-Then-Else Expressions}

The \texttt{if-then-else} expression evaluates a condition and returns one of two values:

\begin{lstlisting}[language=ConstructLang]
if x > 0 then "positive" else "negative"
\end{lstlisting}

The \texttt{else} clause is required. Both branches must return the same type:

\begin{lstlisting}[language=ConstructLang]
let result: String =
  if age >= 18
    then "adult"
    else "minor"
\end{lstlisting}

If-expressions can be nested:

\begin{lstlisting}[language=ConstructLang]
let category: String =
  if x > 0
    then
      if x < 10
        then "small"
        else "large"
    else "negative"
\end{lstlisting}

\section{Match Expressions}

Pattern matching with \texttt{match} provides powerful destructuring:

\begin{lstlisting}[language=ConstructLang]
let desc: String =
  match value:
    | 0 -> "zero"
    | 1 -> "one"
    | _ -> "other"
\end{lstlisting}

The underscore \texttt{\_} matches anything:

\begin{lstlisting}[language=ConstructLang]
let status: String =
  match code:
    | 200 -> "success"
    | 404 -> "not found"
    | 500 -> "server error"
    | _   -> "unknown"
\end{lstlisting}

\chapter{Functions}

\section{Defining Functions}

Functions are defined with \texttt{fn}. The lambda syntax is most common:

\begin{lstlisting}[language=ConstructLang]
# Single parameter
square: Int -> Int = fn (x: Int) -> x * x

# Multiple parameters (curried)
add: Int -> Int -> Int = fn (a: Int) -> fn (b: Int) -> a + b
\end{lstlisting}

\section{Function Application}

Call functions by name with arguments:

\begin{lstlisting}[language=ConstructLang]
let result = square(5)        # result = 25
let sum = add(3)(4)           # sum = 7

# Or use pipe syntax (discussed later)
let sum2 = 3 |> add(4)        # sum2 = 7
\end{lstlisting}

\section{Anonymous Functions}

Create functions without names:

\begin{lstlisting}[language=ConstructLang]
fn (x: Int) -> x * 2
fn (a: Int) -> fn (b: Int) -> a + b
\end{lstlisting}

\section{Function Composition}

Functions can return other functions (higher-order functions):

\begin{lstlisting}[language=ConstructLang]
# make_multiplier creates a multiplier function
make_multiplier: Int -> (Int -> Int) =
  fn (factor: Int) -> fn (x: Int) -> x * factor

# Use it to create specific multipliers
times_two = make_multiplier(2)
times_three = make_multiplier(3)

let result1 = times_two(10)    # result1 = 20
let result2 = times_three(10)  # result2 = 30
\end{lstlisting}

\chapter{Input and Output}

\section{Output Functions}

\subsection{show() - Output Without Newline}

Output one or more values without a trailing newline:

\begin{lstlisting}[language=ConstructLang]
show("Hello ")
show("World")
\end{lstlisting}

\subsection{showln() - Output With Newline}

Output one or more values with a trailing newline:

\begin{lstlisting}[language=ConstructLang]
showln("Hello, Construct!")
showln("Numbers: ", 1, 2, 3)
\end{lstlisting}

The \texttt{show()} and \texttt{showln()} functions work with any type:

\begin{lstlisting}[language=ConstructLang]
let x: Int = 42
let y: Float = 3.14
let name: String = "Alice"

showln("Integer: ", x)
showln("Float: ", y)
showln("String: ", name)
showln("Boolean: ", true)
\end{lstlisting}

\subsection{Variadic Output}

Both \texttt{show()} and \texttt{showln()} support multiple arguments:

\begin{lstlisting}[language=ConstructLang]
show("A", "B", "C")            # ABC
showln("ID: ", 123, " Status: ", "active")
# Outputs: ID: 123 Status: active
\end{lstlisting}

\section{Input Functions}

\subsection{input() - Read Input With Prompt}

Read a line of input with an optional prompt:

\begin{lstlisting}[language=ConstructLang]
let name: String = input("Enter your name: ")
showln("Hello, ", name)
\end{lstlisting}

The prompt is displayed immediately before the input is read.

\subsection{Advanced Input}

Convert input to other types:

\begin{lstlisting}[language=ConstructLang]
let age_str: String = input("Enter your age: ")
let age: Int = toInt(age_str)
showln("You are ", age, " years old")
\end{lstlisting}

\chapter{Standard Library}

\section{String Functions}

\subsection{len() - String Length}

Get the number of characters in a string:

\begin{lstlisting}[language=ConstructLang]
let text: String = "Hello"
let length: Int = len(text)    # length = 5
\end{lstlisting}

\subsection{concat() - Concatenate Strings}

Join two strings:

\begin{lstlisting}[language=ConstructLang]
let greeting: String = concat("Hello", " World")
# greeting = "Hello World"
\end{lstlisting}

Alternatively, use the \texttt{+} operator:

\begin{lstlisting}[language=ConstructLang]
let greeting: String = "Hello" + " " + "World"
\end{lstlisting}

\subsection{upper() - Convert to Uppercase}

Convert a string to uppercase:

\begin{lstlisting}[language=ConstructLang]
let text: String = "hello"
let upper_text: String = upper(text)   # "HELLO"
\end{lstlisting}

\subsection{lower() - Convert to Lowercase}

Convert a string to lowercase:

\begin{lstlisting}[language=ConstructLang]
let text: String = "HELLO"
let lower_text: String = lower(text)   # "hello"
\end{lstlisting}

\subsection{trim() - Remove Whitespace}

Remove leading and trailing whitespace:

\begin{lstlisting}[language=ConstructLang]
let text: String = "  hello  "
let trimmed: String = trim(text)       # "hello"
\end{lstlisting}

\subsection{eq() - Test String Equality}

Test if two strings are equal:

\begin{lstlisting}[language=ConstructLang]
let result: Bool = eq("hello", "hello")   # true
let result2: Bool = eq("hello", "HELLO")  # false
\end{lstlisting}

\subsection{starts\_with() - Check Prefix}

Test if a string starts with another:

\begin{lstlisting}[language=ConstructLang]
let result: Bool = starts_with("hello", "hel")   # true
let result2: Bool = starts_with("hello", "bye")  # false
\end{lstlisting}

\subsection{ends\_with() - Check Suffix}

Test if a string ends with another:

\begin{lstlisting}[language=ConstructLang]
let result: Bool = ends_with("hello", "lo")   # true
let result2: Bool = ends_with("hello", "he")  # false
\end{lstlisting}

\section{Type Conversion Functions}

\subsection{toStr() - Convert to String}

Convert any value to a string:

\begin{lstlisting}[language=ConstructLang]
let s1: String = toStr(42)          # "42"
let s2: String = toStr(3.14)        # "3.14"
let s3: String = toStr(true)        # "true"
\end{lstlisting}

\subsection{toInt() - Convert to Integer}

Parse a string as an integer:

\begin{lstlisting}[language=ConstructLang]
let n: Int = toInt("42")            # 42
let n2: Int = toInt("123")          # 123
\end{lstlisting}

\subsection{toFloat() - Convert to Float}

Parse a string as a float:

\begin{lstlisting}[language=ConstructLang]
let f: Float = toFloat("3.14")      # 3.14
let f2: Float = toFloat("2.5")      # 2.5
\end{lstlisting}

\subsection{toBool() - Convert to Boolean}

Parse a string as a boolean:

\begin{lstlisting}[language=ConstructLang]
let b: Bool = toBool("true")        # true
let b2: Bool = toBool("false")      # false
\end{lstlisting}

\section{Math Functions}

\subsection{abs() - Absolute Value}

Get the absolute value:

\begin{lstlisting}[language=ConstructLang]
let x: Int = abs(-42)               # 42
let y: Float = abs(-3.14)           # 3.14
\end{lstlisting}

\subsection{max() - Maximum}

Get the maximum of two values:

\begin{lstlisting}[language=ConstructLang]
let m1: Int = max(10, 20)           # 20
let m2: Float = max(3.14, 2.71)     # 3.14
\end{lstlisting}

\subsection{min() - Minimum}

Get the minimum of two values:

\begin{lstlisting}[language=ConstructLang]
let m1: Int = min(10, 20)           # 10
let m2: Float = min(3.14, 2.71)     # 2.71
\end{lstlisting}

\subsection{clamp() - Constrain to Range}

Constrain a value within a range:

\begin{lstlisting}[language=ConstructLang]
let c1: Int = clamp(5, 0, 10)       # 5
let c2: Int = clamp(-5, 0, 10)      # 0
let c3: Int = clamp(15, 0, 10)      # 10
\end{lstlisting}

\subsection{round() - Round to Integer}

Round a float to the nearest integer:

\begin{lstlisting}[language=ConstructLang]
let r1: Int = round(3.2)            # 3
let r2: Int = round(3.7)            # 4
\end{lstlisting}

\subsection{floor() - Round Down}

Round down to the nearest integer:

\begin{lstlisting}[language=ConstructLang]
let f1: Int = floor(3.7)            # 3
let f2: Int = floor(3.2)            # 3
\end{lstlisting}

\subsection{ceil() - Round Up}

Round up to the nearest integer:

\begin{lstlisting}[language=ConstructLang]
let c1: Int = ceil(3.2)             # 4
let c2: Int = ceil(3.7)             # 4
\end{lstlisting}

\section{Utility Functions}

\subsection{nl() - Print Newline}

Print a newline:

\begin{lstlisting}[language=ConstructLang]
show("Line 1")
nl()
show("Line 2")
\end{lstlisting}

\subsection{sleep() - Pause Execution}

Pause for a number of seconds:

\begin{lstlisting}[language=ConstructLang]
showln("Waiting...")
sleep(2)
showln("Done!")
\end{lstlisting}

\subsection{clear\_screen() - Clear Terminal}

Clear the terminal screen:

\begin{lstlisting}[language=ConstructLang]
clear_screen()
showln("Fresh start")
\end{lstlisting}

\chapter{Practical Examples}

\section{Example 1: Hello, Name}

A program that greets the user:

\begin{lstlisting}[language=ConstructLang]
let name: String = input("What's your name? ")
showln("Hello, ", name, "!")
\end{lstlisting}

\section{Example 2: Sum of Numbers}

Calculate the sum of user-entered numbers:

\begin{lstlisting}[language=ConstructLang]
let a_str: String = input("Enter first number: ")
let b_str: String = input("Enter second number: ")

let a: Int = toInt(a_str)
let b: Int = toInt(b_str)

let sum: Int = a + b

showln("The sum is: ", sum)
\end{lstlisting}

\section{Example 3: Factorial}

Calculate factorial using recursion:

\begin{lstlisting}[language=ConstructLang]
fn factorial(n: Int): Int {
  if n <= 1
    then 1
    else n * factorial(n - 1)
}

let n: Int = 5
let result: Int = factorial(n)
showln(n, "! = ", result)
\end{lstlisting}

\section{Example 4: String Processing}

Process and display string information:

\begin{lstlisting}[language=ConstructLang]
let text: String = input("Enter some text: ")

showln("Original: ", text)
showln("Length: ", len(text))
showln("Uppercase: ", upper(text))
showln("Lowercase: ", lower(text))
\end{lstlisting}

\section{Example 5: Simple Calculator}

A basic calculator program:

\begin{lstlisting}[language=ConstructLang]
let num1_str: String = input("First number: ")
let operator: String = input("Operator (+, -, *, /): ")
let num2_str: String = input("Second number: ")

let num1: Int = toInt(num1_str)
let num2: Int = toInt(num2_str)

let result: Int =
  if operator == "+"
    then num1 + num2
  else if operator == "-"
    then num1 - num2
  else if operator == "*"
    then num1 * num2
  else if operator == "/"
    then num1 / num2
  else 0

showln("Result: ", result)
\end{lstlisting}

\chapter{Advanced Features}

\section{Higher-Order Functions}

Functions that take or return other functions enable powerful abstractions:

\begin{lstlisting}[language=ConstructLang]
# A function that applies a function to a value twice
apply_twice: (Int -> Int) -> Int -> Int =
  fn (f: Int -> Int) -> fn (x: Int) -> f(f(x))

# Create a function that doubles a number
double: Int -> Int = fn (x: Int) -> x * 2

# Apply double twice to get quadruple
let quadruple = apply_twice(double)
let result = quadruple(5)      # result = 20
\end{lstlisting}

\section{Partial Application}

Apply some arguments to a function to create a new function:

\begin{lstlisting}[language=ConstructLang]
# Original function taking two arguments
add: Int -> Int -> Int =
  fn (a: Int) -> fn (b: Int) -> a + b

# Partially apply to create add_five
add_five: Int -> Int = add(5)

# Use the specialized function
let result = add_five(10)      # result = 15
\end{lstlisting}

\section{Function Composition}

Combine multiple functions into one:

\begin{lstlisting}[language=ConstructLang]
# Define simple functions
increment: Int -> Int = fn (x: Int) -> x + 1
double: Int -> Int = fn (x: Int) -> x * 2
square: Int -> Int = fn (x: Int) -> x * x

# Compose them manually
let result = square(double(increment(5)))  # (5+1)*2)^2 = 144
\end{lstlisting}

\chapter{Common Patterns}

\section{Pattern: Map and Filter}

Process collections with higher-order functions:

\begin{lstlisting}[language=ConstructLang]
# Conceptual: filter values based on a condition
let numbers: String = "1 2 3 4 5"
let large: Bool = if value > 2 then true else false
\end{lstlisting}

\section{Pattern: Recursive Accumulation}

Use recursion to build results:

\begin{lstlisting}[language=ConstructLang]
fn sum_to(n: Int): Int {
  if n <= 0
    then 0
    else n + sum_to(n - 1)
}

let result = sum_to(10)     # 55
\end{lstlisting}

\section{Pattern: Conditional Processing}

Combine if-expressions for powerful control flow:

\begin{lstlisting}[language=ConstructLang]
let grade: String =
  if score >= 90
    then "A"
  else if score >= 80
    then "B"
  else if score >= 70
    then "C"
  else if score >= 60
    then "D"
  else
    "F"
\end{lstlisting}

\chapter{Troubleshooting}

\section{Common Errors}

\subsection{Type Mismatch}

\textbf{Error:} ``Type mismatch: expected Int, got String''

\textbf{Solution:} Use type conversion functions:

\begin{lstlisting}[language=ConstructLang]
let input_str = input("Enter a number: ")
let number: Int = toInt(input_str)   # Convert first!
\end{lstlisting}

\subsection{Undefined Variable}

\textbf{Error:} ``Undefined variable: x''

\textbf{Solution:} Make sure to bind variables before using them:

\begin{lstlisting}[language=ConstructLang]
let x: Int = 42
showln(x)          # OK: x is defined
\end{lstlisting}

\subsection{Indentation Errors}

\textbf{Error:} ``Unexpected token''

\textbf{Solution:} Check indentation (uses spaces, not tabs):

\begin{lstlisting}[language=ConstructLang]
# Correct: consistent indentation
if true
  then showln("yes")
  else showln("no")

# Wrong: inconsistent indentation
if true
then showln("yes")
  else showln("no")
\end{lstlisting}

\appendix

\chapter{Installation and Setup}

\section{Building from Source}

\subsection{Prerequisites}

\begin{itemize}
    \item A C++ compiler (GCC, Clang, or MSVC)
    \item Meson build system
    \item LLVM 15 or later
    \item CMake (optional, for LLVM configuration)
\end{itemize}

\subsection{Build Steps}

\begin{lstlisting}[language=bash]
# Clone or download the Construct repository
cd Construct

# Configure the build
meson setup build

# Compile
meson compile -C build

# Install (optional)
meson install -C build
\end{lstlisting}

\section{Writing Your First Program}

\subsection{Create a File}

Create \texttt{hello.ct}:

\begin{lstlisting}[language=ConstructLang]
showln("Hello, Construct!")
\end{lstlisting}

\subsection{Compile and Run}

\begin{lstlisting}[language=bash]
# Compile to LLVM IR
./build/construct hello.ct

# View the generated IR
cat hello.ll

# Compile to native executable
./build/construct hello.ct -c
clang++ -o hello hello.o ./build/libconstruct_stdlib.so -lm

# Run the program
./hello
\end{lstlisting}

\chapter{Quick Reference}

\section{Keywords}

\begin{center}
\begin{tabular}{|l|l|}
\hline
\textbf{Keyword} & \textbf{Purpose} \\
\hline
\texttt{let} & Bind a variable \\
\texttt{fn} & Define a function \\
\texttt{if} & Conditional expression \\
\texttt{then} & Then branch of if \\
\texttt{else} & Else branch of if \\
\texttt{match} & Pattern matching \\
\texttt{type} & Type alias \\
\texttt{import} & Import module \\
\texttt{true} & Boolean true \\
\texttt{false} & Boolean false \\
\hline
\end{tabular}
\end{center}

\section{Built-in Types}

\begin{center}
\begin{tabular}{|l|l|}
\hline
\textbf{Type} & \textbf{Description} \\
\hline
\texttt{Int} & 32-bit integer \\
\texttt{Float} & 64-bit floating point \\
\texttt{Bool} & Boolean \\
\texttt{String} & Text string \\
\hline
\end{tabular}
\end{center}

\section{I/O Functions}

\begin{center}
\begin{tabular}{|l|l|}
\hline
\textbf{Function} & \textbf{Purpose} \\
\hline
\texttt{show(...)} & Output without newline \\
\texttt{showln(...)} & Output with newline \\
\texttt{input(prompt)} & Read input with prompt \\
\texttt{nl()} & Print newline \\
\hline
\end{tabular}
\end{center}

\section{String Functions}

\begin{center}
\begin{tabular}{|l|l|}
\hline
\textbf{Function} & \textbf{Purpose} \\
\hline
\texttt{len(s)} & String length \\
\texttt{concat(a, b)} & Concatenate strings \\
\texttt{upper(s)} & Convert to uppercase \\
\texttt{lower(s)} & Convert to lowercase \\
\texttt{trim(s)} & Remove whitespace \\
\texttt{eq(a, b)} & Test equality \\
\texttt{starts\_with(s, p)} & Check prefix \\
\texttt{ends\_with(s, p)} & Check suffix \\
\hline
\end{tabular}
\end{center}

\section{Math Functions}

\begin{center}
\begin{tabular}{|l|l|}
\hline
\textbf{Function} & \textbf{Purpose} \\
\hline
\texttt{abs(x)} & Absolute value \\
\texttt{max(a, b)} & Maximum \\
\texttt{min(a, b)} & Minimum \\
\texttt{clamp(x, min, max)} & Constrain range \\
\texttt{round(x)} & Round to nearest \\
\texttt{floor(x)} & Round down \\
\texttt{ceil(x)} & Round up \\
\hline
\end{tabular}
\end{center}

\chapter{Solutions to Exercises}

\section{Exercise 1: Temperature Converter}

Convert Fahrenheit to Celsius:

\begin{lstlisting}[language=ConstructLang]
let f_str: String = input("Enter temperature in F: ")
let f: Float = toFloat(f_str)
let c: Float = (f - 32.0) * (5.0 / 9.0)
showln(toStr(f), "F = ", toStr(c), "C")
\end{lstlisting}

\section{Exercise 2: Password Validator}

Check if a password is strong:

\begin{lstlisting}[language=ConstructLang]
let password: String = input("Enter password: ")
let is_long: Bool = len(password) >= 8

let message: String =
  if is_long
    then "Password is strong"
    else "Password is too short"

showln(message)
\end{lstlisting}

\section{Exercise 3: Number Guessing Game}

\begin{lstlisting}[language=ConstructLang]
let secret: Int = 42
let guess_str: String = input("Guess the number: ")
let guess: Int = toInt(guess_str)

let message: String =
  if guess == secret
    then "Correct!"
  else if guess < secret
    then "Too low"
  else
    "Too high"

showln(message)
\end{lstlisting}

\chapter{Resources}

\section{Online References}

\begin{itemize}
    \item \textbf{Construct GitHub}: \url{https://github.com/charlie-sans/construct}
    \item \textbf{LLVM Documentation}: \url{https://llvm.org/docs/}
    \item \textbf{Functional Programming}: \url{https://en.wikipedia.org/wiki/Functional_programming}
\end{itemize}

\section{Recommended Reading}

\begin{itemize}
    \item \textit{Composing Software} by Eric Elliott
    \item \textit{Structure and Interpretation of Computer Programs} by Abelson \& Sussman
    \item \textit{Learn You a Haskell} by Lipovača
\end{itemize}

\backmatter

\chapter*{Conclusion}

You now have a comprehensive understanding of the Construct programming language. Whether you're writing simple scripts or complex programs, you have the tools and knowledge to express your ideas clearly and effectively in Construct.

Remember these key principles:

\begin{enumerate}
    \item \textbf{Everything is an expression}: Think in terms of values and transformations
    \item \textbf{Functions are first-class}: Combine functions to build powerful abstractions
    \item \textbf{Types help you}: Use them to catch errors early and clarify intent
    \item \textbf{Immutability simplifies reasoning}: Avoid state changes when possible
    \item \textbf{Compose building blocks}: Create complex programs from simple, reusable pieces
\end{enumerate}

Continue exploring, experimenting, and building. The Construct community is here to support you.

Happy coding!

\end{document}
